\subsection{截面簇与映射簇}

15.3.1（“截面空间上的 Banach 空间结构”）。设 B 是一个 \(C^{r}\) 类紧微分流形；设 M 是一个以 B 为底、\(C^{r}\) 类的 K 上的向量丛（7.3.4）；用 \(M_{r}\) 表示 M 的底层 \(C^{r}\) 类微分流形，并用 \(\mathcal{S}_{\mathrm{M}}^{\tau}(\mathrm{B})\) 表示 M 的 \(C^{r}\) 类截面所成的 K-向量空间（7.4.1）。在 \(\mathcal{S}_{\mathrm{M}}^{\tau}(\mathrm{B})\) 上，由 \(C^{r}\) -一致收敛拓扑（在 \(\mathcal{C}^{\tau}(\mathrm{B};\mathrm{M}_{r})\) 上，12.3.10）诱导的拓扑使 \(\mathcal{S}_{\mathrm{M}}^{\tau}(\mathrm{B})\) 成为一个 K-Banach 空间。

设 \((\mathrm{U}_{i})_{i\in\mathbb{I}}\) 是 B 的一个有限开覆盖，并且对每个 \(i\in I\) ，设 \(c_{i}=(\mathrm{U}_{i},\varphi_{i},\mathrm{E}_{i})\) 是 B 的一个卡，\(t_{i}=(\mathrm{U}_{i},\psi_{i},\mathrm{F}_{i})\) 是 M 的一个 K-向量卡（8.8.1）。设每个 \(E_{i}\) （相应地，\(F_{i}\) ）都赋有与它在 R（相应地，K）上的 Banach 空间结构相容的范数，并且 \((\mathrm{V}_{i})_{i\in\mathbb{I}}\) 是 B 的一个闭覆盖，使得 \(V_{i}\subset U_{i}\) 对所有 \(i\in I\) 成立。对 \(f\in\mathcal{S}_{\mathrm{M}}^{\tau}(\mathrm{B})\) 与 \(i\in I\) ，用 \(f_{i}\) 表示由 \(\varphi_{i}(\mathrm{U}_{i})\) 到 \(F_{i}\) 的映射，它由下述复合得到：

\[
\varphi_ {i} \left(\mathrm{U} _ {i}\right) \xrightarrow {\varphi_ {i} ^ {- 1}} \mathrm{U} _ {i} \xrightarrow {f} \mathrm{M} \left| \mathrm{U} _ {i} \xrightarrow {\psi_ {i}} \mathrm{U} _ {i} \times \mathrm{F} _ {i} \xrightarrow {\operatorname{pr} _ {2}} \mathrm{F} _ {i}. \right.
\]

规定

\[
\| f \| = \operatorname{Sup} _ {i \in I, p \leqslant r, x \in V _ {i}} \| D ^ {p} f _ {i} (\varphi_ {i} (x)) \|.
\]

映射 \(f\mapsto\|f\|\) 是 \(\mathcal{S}_{\mathrm{M}}^{\tau}(\mathrm{B})\) 上的一个范数，它给 \(\mathcal{S}_{\mathrm{M}}^{\tau}(\mathrm{B})\) 一个 K 上的 Banach 空间结构。

若 N 是空间 M 的一个开子集，则集合 \(\mathcal{S}^{\tau}(\mathbf{B};\mathbf{N})\) （由截面 \(f\in\mathcal{S}_{\mathbf{M}}^{\tau}(\mathbf{B})\) 中满足 \(f(\mathbf{B})\subset\mathbf{N}\) 者组成）在 \(\mathcal{S}_{\mathbf{M}}^{\tau}(\mathbf{B})\) 中是开的。

15.3.2. 设 B 是一个 \(C^{r}\) 类紧微分流形，X 是 B 上的一个 \(C_{K}^{r,s}\) 类（ \(s \geqslant r + 1\) ）流形。若 N 是 X 的一个开子集，则用 \(\mathcal{S}^{r}(B; N)\) 表示 B 上的像含于 N 的 \(C^{r}\) 类截面所成的集合。在 \(\mathcal{S}^{r}(B; X)\) 上存在、且仅存在一个 \(C^{s-r}\) 类 K-流形结构，使得对每个截面 \(\sigma \in \mathcal{S}^{r}(B; X)\) 以及每个管状邻域 (M, N, \(\varphi\) )（\(C_{K}^{r,s}\) 类，\(\sigma\) 的，15.2.7），三元组 ( \(\mathcal{S}^{r}(B; N)\) , \(\varphi^{*}\) , \(\mathcal{S}_{M}^{r}(B)\) )——其中 \(\varphi^{*}\) 是映射 \(f \mapsto \varphi \circ f\) ——是 K-流形 \(\mathcal{S}^{r}(B; X)\) 的一个卡。

在下文中，我们给 \(\mathcal{S}^{r}(\mathbf{B};\mathbf{X})\) 赋以这一结构。

15.3.3. 设 X 是一个 \(C^{r}\) 类紧微分流形，Y 是一个 \(C^{s}(s \geqslant r + 1)\) 类 K-流形。设 \(\mathcal{C}^{r}(X; Y)\) 是从 X 到 Y 中的 \(C^{r}\) 类映射所成的集合（换句话说，取值于 Y 的底层 \(C^{r}\) 类流形中，参见 5.13.1 与 5.14.2）。我们给 \(\mathcal{C}^{r}(X; Y)\) 赋以一个 \(C^{s-r}\) 类 K-簇结构，它由与 \(\mathcal{S}^{r}(X; X \times Y)\) 的等同得到（参见 15.2.3, (ii) 与 15.3.2）。称它为从 X 到 Y 中的 \(C^{r}\) 类映射簇。它的拓扑是 \(C^{r}\) -一致收敛拓扑（12.3.10）。

假设 Y 是分离的；则 \(\mathcal{C}^{r}(X;Y)\) 也是分离的。此外，从 X 到 Y 中的 \(C^{r}\) 类浸入（相应地，嵌入、浸没、满射浸没、平展态射、同构）所成的集合是 \(\mathcal{C}^{r}(X;Y)\) 的一个开子集。

设 \(X'\) 是一个 \(C^r\) 类紧微分流形，\(\varphi: X' \to X\) 是 \(C^r\) 类微分流形的一个态射，而 \(\psi: Y \to Y'\) 是 \(C^s\) 类 K-簇的一个态射。映射 \(f \mapsto \psi \circ f \circ \varphi\) 是一个 \(C^{s-r}\) 类态射（从 \(\mathscr{C}^r(X; Y)\) 到 \(\mathscr{C}^r(X'; Y')\) 中）。

15.3.4（“结构弱化”）。保留 15.3.3 的假设与记号。此外，设 \(s' \in \mathbf{N}_{\mathbf{R}}\) 满足 \(r < s' \leqslant s\) ，并设 \(\mathbf{Y}_{s'}\) 是一个 \(\mathbf{C}^{s'}\) 类实流形，它以 \(\mathbf{Y}\) 为底层（5.13 与 5.14.2）。\(\mathbf{C}^{s'-r}\) 类实流形结构（在 \(\mathcal{C}^r(\mathbf{X}; \mathbf{Y}_{s'})\) 上）是 \(\mathbf{C}^{s-r}\) 类 K-流形结构（在 \(\mathcal{C}^r(\mathbf{X}; \mathbf{Y})\) 上）的底层。

15.3.5. 设 X 与 X' 分别是 \(C^{r}\) 类与 \(C^{r'}\) 类的紧微分流形，且 \(r' < r\) ；设 Y 是一个 \(C^{s}\) 类 K-簇，\(s \geqslant r + 1\) 。令 \(t = \operatorname{Inf}(r - r', s - r)\) 。映射 \((f, g) \mapsto g \circ f\) 是一个 \(C^{t}\) 类态射（从 \(\mathcal{C}^{r}(X'; X) \times \mathcal{C}^{r}(X; Y)\) 到 \(\mathcal{C}^{r'}(X'; Y)\) 中）。

15.3.6（“\(\mathcal{C}^r (\mathrm{X};\mathrm{Y})\) 的切空间”）。在 15.3.3 的假设之下，设 \(f\in \mathcal{C}^r (\mathrm{X};\mathrm{Y})\) ，并设 \(\xi\) 是 \(f\) 处切于流形 \(\mathcal{C}^r (\mathrm{X};\mathrm{Y})\) 的一个切向量。若 \(x\in \mathrm{X}\) ，则用 \(\varepsilon_{x}\) 表示映射 \(g\mapsto g(x)\) （从 \(\mathcal{C}^r (\mathrm{X};\mathrm{Y})\) 到 Y 中）；它是一个 \(\mathbf{C}^{s - r}\) 类态射。\(\xi\) 在 \(\mathrm{T}_f(\varepsilon_x)\) 下的像是一个元素 \(\xi_{x}\) ，它属于 \(\mathrm{T}_{f(x)}(\mathrm{Y})\) ；映射 \(x\mapsto \xi_{x}\) 是一个 \(\mathbf{C}^r\) 类提升，即 \(f\) 在 T(Y) 中的提升（参见 8.6.1），并且可以把它与一个 \(\mathbf{C}^r\) 类截面（属于以 X 为底的向量丛 \(f^{*}\mathrm{T}(\mathrm{Y})\)）等同起来。这样，就得到了切空间 \(\mathrm{T}_f(\mathcal{C}^r (\mathrm{X};\mathrm{Y}))\) 到 Banach 空间 \(\mathcal{S}_{f^{*}\mathrm{T}(\mathrm{Y})}^r (\mathrm{X})\) 上的一个同构。

15.3.7（“\(\mathcal{C}^{p}(\mathrm{X};\mathcal{C}^{q}(\mathrm{Y};\mathrm{Z}))\) 的解释”）。设 X 与 Y 分别是 \(\mathbf{C}^{p}\) 类与 \(\mathbf{C}^{q}\) 类的紧微分流形（ \(1\leqslant p\leqslant+\infty,1\leqslant q<+\infty\) ），并设 Z 是一个 \(\mathbf{C}^{s}\) 类 K-簇，\(s\geqslant p+q\) 。设 \(f\colon\mathrm{X}\times\mathrm{Y}\to\mathrm{Z}\) 是一个连续映射；

对每个 \(x \in X\) ，用 \(f_x\) 表示从 Y 到 Z 的映射 \(y \mapsto f(x, y)\) 。下列两个性质等价：

\begin{enumerate}
\def\labelenumi{\alph{enumi})}
\item
  对任意卡 \(c = (\mathrm{U}, \varphi, \mathrm{E})\) （\(\mathbf{X}\) 的）、任意卡 \(d = (\mathrm{V}, \psi, \mathrm{F})\) （\(\mathbf{Y}\) 的）与任意卡 \(e = (\mathrm{W}, \theta, \mathrm{G})\) （\(\mathbf{Z}\) 的），只要 \(f(\mathrm{U} \times \mathrm{V}) \subset \mathrm{W}\) ，偏导数 \(\mathrm{D}_{\mathbb{E}}^{p'} \mathrm{D}_{\mathbb{F}}^{q'} f'\) （其中 \(f'\) 是 \(f\) 在卡 \(c \times d\) 与 \(e\) 中的表达式，5.3.2）对 \(p' \leqslant p\) 与 \(q' \leqslant q\) 都存在且连续。
\item
  对每个 \(x \in X\) ，映射 \(f_{x}\) 是 \(C^{q}\) 类的，且映射 \(x \mapsto f_{x}\) 是一个 \(C^{p}\) 类态射（从 X 到簇 \(\mathcal{C}^{q}(Y; Z)\) 中）。
\end{enumerate}

\pandocbounded{\includegraphics[keepaspectratio]{images/part001_50717dae.jpg}}

15.3.8. 设 Y 是一个 \(C^{s}\) 类紧 K-簇，\(s \geqslant r + 1\) ，并设 \(\text{Diff}^{r}(Y)\) 是 Y 的底层 \(C^{r}\) 类实流形 \(Y_{r}\) 的 \(C^{r}\) 类自同构所成的群。集合 \(\text{Diff}^{r}(Y)\) 在 \(\mathcal{C}^{r}(Y; Y) = \mathcal{C}^{r}(Y_{r}; Y)\) 中是开的；我们给它赋以一个 \(C^{s-r}\) 类 K-簇结构（由 \(\mathcal{C}^{r}(Y; Y)\) 诱导）。复合律 \((f, g) \mapsto g \circ f\) 使 \(\text{Diff}^{r}(Y)\) 成为一个拓扑群。另一方面，\(C^{s-r}\) 类流形结构一般说来与 \(\text{Diff}^{r}(Y)\) 的群结构不相容。对 \(f \in \text{Diff}^{r}(Y)\) ，映射 \(g \mapsto g \circ f\) （从 \(\text{Diff}^{r}(Y)\) 到其自身）是 \(C^{s-r}\) 类的，但映射 \(g \mapsto f \circ g\) 一般说来不是 \(C^{1}\) 类的。特别地，\(\text{Diff}^{r}(Y)\) 一般说来不是一个李群。\footnote{假设 \(s = \infty\)。设 \(\text{Diff}^{\infty}(Y)\) 是 Y 的 \(C^{\infty}\) 类自同构所成的群；给它赋以使诸典范单射 \(\text{Diff}^{\infty}(Y) \to \text{Diff}^{r}(Y)\)（对每个整数 \(r \geqslant 1\)）都连续的最粗拓扑。另一方面，设 \(\text{Diff}^{0}(Y)\) 是 Y 的同胚所成的群，赋以一致收敛拓扑（TG, X, § 3, no. 5, 命题 11）。若 \(\dim Y \geqslant 1\)，则不可能找到一个李群 G 以及从 \(\text{Diff}^{\infty}(Y)\) 到 G 和从 G 到 \(\text{Diff}^{0}(Y)\) 的连续同态，使其复合恰为典范单射（参见 LIE, III, § 4, 习题 7）。}

